\section{JOINT DECLARATION ON THE DOCTRINE OF
JUSTIFICATION}\label{joint-declaration-on-the-doctrine-of-justification}

\subsection{\texorpdfstring{\emph{by the Lutheran World Federation and
the Catholic
Church}}{by the Lutheran World Federation and the Catholic Church}}\label{by-the-lutheran-world-federation-and-the-catholic-church}

\section{~}\label{section}

\subsubsection{Preamble}\label{preamble}

1.~ ~The doctrine of justification was of central importance for the
Lutheran Reformation of the sixteenth century. It was held to be the
"first and chief
article"\phantomsection\label{r1}{}\hyperref[1]{{[}1{]}}
\textsuperscript{} and at the same time the "ruler and judge over all
other Christian
doctrines."\phantomsection\label{r2}{}\hyperref[2]{{[}2{]}} The doctrine
of justification was particularly asserted and defended in its
Reformation shape and special valuation over against the Roman Catholic
Church and theology of that time, which in turn asserted and defended a
doctrine of justification of a different character. From the Reformation
perspective, justification was the crux of all the disputes. Doctrinal
condemnations were put forward both in the Lutheran
Confessions\phantomsection\label{r3}{}\hyperref[3]{{[}3{]}} and by the
Roman Catholic Church\textquotesingle s Council of Trent. These
condemnations are still valid today and thus have a church-dividing
effect.

2.~ ~For the Lutheran tradition, the doctrine of justification has
retained its special status. Consequently it has also from the beginning
occupied an important place in the official Lutheran-Roman Catholic
dialogue.

3.~ ~Special attention should be drawn to the following reports: "The
Gospel and the Church"
(1972)\phantomsection\label{r4}{}\hyperref[4]{{[}4{]}} and "Church and
Justification" (1994)\phantomsection\label{r5}{}\hyperref[5]{{[}5{]}} by
the Lutheran-Roman Catholic Joint Commission, "Justification by Faith"
(1983)\phantomsection\label{r6}{}\hyperref[6]{{[}6{]}} of the
Lutheran-Roman Catholic dialogue in the USA and "The Condemnations of
the Reformation Era - Do They Still Divide?"
(1986)\phantomsection\label{r7}{}\hyperref[7]{{[}7{]}} by the Ecumenical
Working Group of Protestant and Catholic theologians in Germany. Some of
these dialogue reports have been officially received by the churches. An
important example of such reception is the binding response of the
United Evangelical-Lutheran Church of Germany to the "Condemnations"
study, made in 1994 at the highest possible level of ecclesiastical
recognition together with the other churches of the Evangelical Church
in Germany.\phantomsection\label{r8}{}\hyperref[8]{{[}8{]}}

4.~ ~In their discussion of the doctrine of justification, all the
dialogue reports as well as the responses show a high degree of
agreement in their approaches and conclusions. The time has therefore
come to take stock and to summarize the results of the dialogues on
justification so that our churches may be informed about the overall
results of this dialogue with the necessary accuracy and brevity, and
thereby be enabled to make binding decisions.

5.~ The present Joint Declaration has this intention: namely, to show
that on the basis of their dialogue the subscribing Lutheran churches
and the Roman Catholic
Church\phantomsection\label{r9}{}\hyperref[9]{{[}9{]}} are now able to
articulate a common understanding of our justification by
God\textquotesingle s grace through faith in Christ. It does not cover
all that either church teaches about justification; it does encompass a
consensus on basic truths of the doctrine of justification and shows
that the remaining differences in its explication are no longer the
occasion for doctrinal condemnations.

6.~ Our Declaration is not a new, independent presentation alongside the
dialogue reports and documents to date, let alone a replacement of them.
Rather, as the appendix of sources shows, it makes repeated reference to
them and their arguments.

7.~ ~Like the dialogues themselves, this Joint Declaration rests on the
conviction that in overcoming the earlier controversial questions and
doctrinal condemnations, the churches neither take the condemnations
lightly nor do they disavow their own past. On the contrary, this
Declaration is shaped by the conviction that in their respective
histories our churches have come to new insights. Developments have
taken place which not only make possible, but also require the churches
to examine the divisive questions and condemnations and see them in a
new light.

~

\subsubsection{1. Biblical Message of
Justification}\label{biblical-message-of-justification}

8.~ ~Our common way of listening to the word of God in Scripture has led
to such new insights. Together we hear the gospel that "God so loved the
world that he gave his only Son, so that everyone who believes in him
may not perish but may have eternal life" (\emph{Jn} 3:16). This good
news is set forth in Holy Scripture in various ways. In the Old
Testament we listen to God\textquotesingle s word about human sinfulness
(\emph{Ps} 51:1-5; \emph{Dan} 9:5f; \emph{Eccl/Qo} 8:9f; \emph{Ezra}
9:6f) and human disobedience (\emph{Gen} 3:1-19; \emph{Neh} 9:16f,26) as
well as of God\textquotesingle s "righteousness" (\emph{Isa} 46:13;
51:5-8; 56:1 {[}cf. 53:11{]}; \emph{Jer} 9:24) and "judgment"
(\emph{Eccl/Qo} 12:14; \emph{Ps} 9:5f; 76:7-9).

9.~ ~In the New Testament diverse treatments of "righteousness" and
"justification" are found in the writings of Matthew (5:10; 6:33;
21:32), John (16:8-11), Hebrews (5:3; 10:37f), and James
(2:14-26).\phantomsection\label{r10}{}\hyperref[10]{{[}10{]}} In
Paul\textquotesingle s letters also, the gift of salvation is described
in various ways, among others: "for freedom Christ has set us free"
(\emph{Gal} 5:1-13; cf. \emph{Rom} 6:7), "reconciled to God" (2
\emph{Cor} 5:18-21; cf. \emph{Rom} 5:11), "peace with God" (\emph{Rom}
5:1), "new creation" (2 \emph{Cor} 5:17), "alive to God in Christ Jesus"
(\emph{Rom} 6:11,23), or "sanctified in Christ Jesus" (cf. 1 \emph{Cor}
1:2; 1:30; 2 \emph{Cor} 1:1). Chief among these is the "justification"
of sinful human beings by God\textquotesingle s grace through faith
(\emph{Rom} 3:23-25), which came into particular prominence in the
Reformation period.

10.~ ~Paul sets forth the gospel as the power of God for salvation of
the person who has fallen under the power of sin, as the message that
proclaims that "the righteousness of God is revealed through faith for
faith" (\emph{Rom} 1:16f) and that grants "justification" (\emph{Rom}
3:21-31). He proclaims Christ as "our righteousness" (1 \emph{Cor}
1:30), applying to the risen Lord what Jeremiah proclaimed about God
himself (\emph{Jer} 23:6). In Christ\textquotesingle s death and
resurrection all dimensions of his saving work have their roots for he
is "our Lord, who was put to death for our trespasses and raised for our
justification" (\emph{Rom} 4:25). All human beings are in need of
God\textquotesingle s righteousness, "since all have sinned and fall
short of the glory of God" (\emph{Rom} 3:23; cf. \emph{Rom} 1:18-3:20;
11:32; \emph{Gal} 3:22). In Galatians (3:6) and Romans (4:3-9), Paul
understands Abraham\textquotesingle s faith (\emph{Gen} 15:6) as faith
in the God who justifies the sinner (\emph{Rom} 4:5) and calls upon the
testimony of the Old Testament to undergird his gospel that this
righteousness will be reckoned to all who, like Abraham, trust in
God\textquotesingle s promise. "For the righteous will live by faith
(Hab 2:4; cf. \emph{Gal} 3:11; \emph{Rom} 1:17). In
Paul\textquotesingle s letters, God\textquotesingle s righteousness is
also God\textquotesingle s power for those who have faith (\emph{Rom}
1:16f; 2 \emph{Cor} 5:21). In Christ he makes it our righteousness (2
\emph{Cor} 5:21). Justification becomes ours through Christ Jesus "whom
God put forward as a sacrifice of atonement by his blood, effective
through faith" (\emph{Rom} 3:25; see 3:21-28). "For by grace you have
been saved through faith, and this is not your own doing; it is the gift
of God - not the result of works" (\emph{Eph} 2:8f).

11.~ ~Justification is the forgiveness of sins (cf. \emph{Rom} 3:23-25;
\emph{Acts} 13:39; \emph{Lk} 18:14), liberation from the dominating
power of sin and death (\emph{Rom} 5:12-21) and from the curse of the
law (\emph{Gal} 3:10-14). It is acceptance into communion with God:
already now, but then fully in God\textquotesingle s coming kingdom
(\emph{Rom} 5:1f). It unites with Christ and with his death and
resurrection (\emph{Rom} 6:5). It occurs in the reception of the Holy
Spirit in baptism and incorporation into the one body (\emph{Rom} 8:1f,
9f; I \emph{Cor} 12:12f). All this is from God alone, for
Christ\textquotesingle s sake, by grace, through faith in "the gospel of
God\textquotesingle s Son" (\emph{Rom} 1:1-3).

12.~ ~The justified live by faith that comes from the Word of Christ
(\emph{Rom} 10:17) and is active through love (\emph{Gal} 5:6), the
fruit of the Spirit (\emph{Gal} 5:22f). But since the justified are
assailed from within and without by powers and desires (\emph{Rom}
8:35-39; \emph{Gal} 5:16-21) and fall into sin (1 \emph{Jn} 1:8,10),
they must constantly hear God\textquotesingle s promises anew, confess
their sins (1 \emph{Jn} 1:9), participate in Christ\textquotesingle s
body and blood, and be exhorted to live righteously in accord with the
will of God. That is why the Apostle says to the justified: "Work out
your own salvation with fear and trembling; for it is God who is at work
in you, enabling you both to will and to work for his good pleasure"
(\emph{Phil} 2:12f). But the good news remains: "there is now no
condemnation for those who are in Christ Jesus" (\emph{Rom} 8:1), and in
whom Christ lives (\emph{Gal} 2:20). Christ\textquotesingle s "act of
righteousness leads to justification and life for all" (\emph{Rom}
5:18).

~

\subsubsection{2. The Doctrine of Justification as Ecumenical
Problem}\label{the-doctrine-of-justification-as-ecumenical-problem}

13.~ ~Opposing interpretations and applications of the biblical message
of justification were in the sixteenth century a principal cause of the
division of the Western church and led as well to doctrinal
condemnations. A common understanding of justification is therefore
fundamental and indispensable to overcoming that division. By
appropriating insights of recent biblical studies and drawing on modern
investigations of the history of theology and dogma, the post-Vatican II
ecumenical dialogue has led to a notable convergence concerning
justification, with the result that this Joint Declaration is able to
formulate a consensus on basic truths concerning the doctrine of
justification. In light of this consensus, the corresponding doctrinal
condemnations of the sixteenth century do not apply to
today\textquotesingle s partner.

~

\subsubsection{3. The Common Understanding of
Justification}\label{the-common-understanding-of-justification}

14.~ ~The Lutheran churches and the Roman Catholic Church have together
listened to the good news proclaimed in Holy Scripture. This common
listening, together with the theological conversations of recent years,
has led to a shared understanding of justification. This encompasses a
consensus in the basic truths; the differing explications in particular
statements are compatible with it.

15.~ ~In faith we together hold the conviction that justification is the
work of the triune God. The Father sent his Son into the world to save
sinners. The foundation and presupposition of justification is the
incarnation, death, and resurrection of Christ. Justification thus means
that Christ himself is our righteousness, in which we share through the
Holy Spirit in accord with the will of the Father. Together we confess:
By grace alone, in faith in Christ\textquotesingle s saving work and not
because of any merit on our part, we are accepted by God and receive the
Holy Spirit, who renews our hearts while equipping and calling us to
good works.\phantomsection\label{r11}{}\hyperref[11]{{[}11{]}}

16.~ ~All people are called by God to salvation in Christ. Through
Christ alone are we justified, when we receive this salvation in faith.
Faith is itself God\textquotesingle s gift through the Holy Spirit who
works through word and sacrament in the community of believers and who,
at the same time, leads believers into that renewal of life which God
will bring to completion in eternal life.

17.~ ~We also share the conviction that the message of justification
directs us in a special way towards the heart of the New Testament
witness to God\textquotesingle s saving action in Christ: it tells us
that as sinners our new life is solely due to the forgiving and renewing
mercy that God imparts as a gift and we receive in faith, and never can
merit in any way.

18.~ ~Therefore the doctrine of justification, which takes up this
message and explicates it, is more than just one part of Christian
doctrine. It stands in an essential relation to all truths of faith,
which are to be seen as internally related to each other. It is an
indispensable criterion which constantly serves to orient all the
teaching and practice of our churches to Christ. When Lutherans
emphasize the unique significance of this criterion, they do not deny
the interrelation and significance of all truths of faith. When
Catholics see themselves as bound by several criteria, they do not deny
the special function of the message of justification. Lutherans and
Catholics share the goal of confessing Christ in all things, who alone
is to be trusted above all things as the one Mediator (1 \emph{Tim}
2:5f) through whom God in the Holy Spirit gives himself and pours out
his renewing gifts. {[}cf. Sources for section 3{]}.

~

\subsubsection{4. Explicating the Common Understanding of
Justification}\label{explicating-the-common-understanding-of-justification}

\paragraph{\texorpdfstring{\emph{4.1 Human Powerlessness and Sin in
Relation to
Justification}}{4.1 Human Powerlessness and Sin in Relation to Justification}}\label{human-powerlessness-and-sin-in-relation-to-justification}

19.~ ~We confess together that all persons depend completely on the
saving grace of God for their salvation. The freedom they possess in
relation to persons and the things of this world is no freedom in
relation to salvation, for as sinners they stand under
God\textquotesingle s judgment and are incapable of turning by
themselves to God to seek deliverance, of meriting their justification
before God, or of attaining salvation by their own abilities.
Justification takes place solely by God\textquotesingle s grace. Because
Catholics and Lutherans confess this together, it is true to say:

20.~ ~When Catholics say that persons "cooperate" in preparing for and
accepting justification by consenting to God\textquotesingle s
justifying action, they see such personal consent as itself an effect of
grace, not as an action arising from innate human abilities.

21.~ ~According to Lutheran teaching, human beings are incapable of
cooperating in their salvation, because as sinners they actively oppose
God and his saving action. Lutherans do not deny that a person can
reject the working of grace. When they emphasize that a person can only
receive (mere passive) justification, they mean thereby to exclude any
possibility of contributing to one\textquotesingle s own justification,
but do not deny that believers are fully involved personally in their
faith, which is effected by God\textquotesingle s Word. {[}cf. Sources
for 4.1{]}.

~

\paragraph{\texorpdfstring{\emph{4.2 Justification as Forgiveness of
Sins and Making
Righteous}}{4.2 Justification as Forgiveness of Sins and Making Righteous}}\label{justification-as-forgiveness-of-sins-and-making-righteous}

22.~ ~We confess together that God forgives sin by grace and at the same
time frees human beings from sin\textquotesingle s enslaving power and
imparts the gift of new life in Christ. When persons come by faith to
share in Christ, God no longer imputes to them their sin and through the
Holy Spirit effects in them an active love. These two aspects of
God\textquotesingle s gracious action are not to be separated, for
persons are by faith united with Christ, who in his person is our
righteousness (1 \emph{Cor} 1:30): both the forgiveness of sin and the
saving presence of God himself. Because Catholics and Lutherans confess
this together, it is true to say that:

23.~ ~When Lutherans emphasize that the righteousness of Christ is our
righteousness, their intention is above all to insist that the sinner is
granted righteousness before God in Christ through the declaration of
forgiveness and that only in union with Christ is one\textquotesingle s
life renewed. When they stress that God\textquotesingle s grace is
forgiving love ("the favor of
God"\phantomsection\label{r12}{}\hyperref[12]{{[}12{]}}), they do not
thereby deny the renewal of the Christian\textquotesingle s life. They
intend rather to express that justification remains free from human
cooperation and is not dependent on the life-renewing effects of grace
in human beings.

24.~ ~When Catholics emphasize the renewal of the interior person
through the reception of grace imparted as a gift to the
believer,\phantomsection\label{r13}{}\hyperref[13]{{[}13{]}} they wish
to insist that God\textquotesingle s forgiving grace always brings with
it a gift of new life, which in the Holy Spirit becomes effective in
active love. They do not thereby deny that God\textquotesingle s gift of
grace in justification remains independent of human cooperation. {[}cf.
Sources for section 4.2{]}.

~

\paragraph{\texorpdfstring{\emph{4.3 Justification by Faith and through
Grace}}{4.3 Justification by Faith and through Grace}}\label{justification-by-faith-and-through-grace}

25.~ ~We confess together that sinners are justified by faith in the
saving action of God in Christ. By the action of the Holy Spirit in
baptism, they are granted the gift of salvation, which lays the basis
for the whole Christian life. They place their trust in
God\textquotesingle s gracious promise by justifying faith, which
includes hope in God and love for him. Such a faith is active in love
and thus the Christian cannot and should not remain without works. But
whatever in the justified precedes or follows the free gift of faith is
neither the basis of justification nor merits it.

26.~ ~According to Lutheran understanding, God justifies sinners in
faith alone (sola fide). In faith they place their trust wholly in their
Creator and Redeemer and thus live in communion with him. God himself
effects faith as he brings forth such trust by his creative word.
Because God\textquotesingle s act is a new creation, it affects all
dimensions of the person and leads to a life in hope and love. In the
doctrine of "justification by faith alone," a distinction but not a
separation is made between justification itself and the renewal of
one\textquotesingle s way of life that necessarily follows from
justification and without which faith does not exist. Thereby the basis
is indicated from which the renewal of life proceeds, for it comes forth
from the love of God imparted to the person in justification.
Justification and renewal are joined in Christ, who is present in faith.

27.~ ~The Catholic understanding also sees faith as fundamental in
justification. For without faith, no justification can take place.
Persons are justified through baptism as hearers of the word and
believers in it. The justification of sinners is forgiveness of sins and
being made righteous by justifying grace, which makes us children of
God. In justification the righteous receive from Christ faith, hope, and
love and are thereby taken into communion with
him.\phantomsection\label{r14}{}\hyperref[14]{{[}14{]}} This new
personal relation to God is grounded totally on God\textquotesingle s
graciousness and remains constantly dependent on the salvific and
creative working of this gracious God, who remains true to himself, so
that one can rely upon him. Thus justifying grace never becomes a human
possession to which one could appeal over against God. While Catholic
teaching emphasizes the renewal of life by justifying grace, this
renewal in faith, hope, and love is always dependent on
God\textquotesingle s unfathomable grace and contributes nothing to
justification about which one could boast before God (\emph{Rom} 3:27).
{[}See Sources for section 4.3{]}.

~

\paragraph{\texorpdfstring{\emph{4.4 The Justified as
Sinner}}{4.4 The Justified as Sinner}}\label{the-justified-as-sinner}

28.~ ~We confess together that in baptism the Holy Spirit unites one
with Christ, justifies, and truly renews the person. But the justified
must all through life constantly look to God\textquotesingle s
unconditional justifying grace. They also are continuously exposed to
the power of sin still pressing its attacks (cf. \emph{Rom} 6:12-14) and
are not exempt from a lifelong struggle against the contradiction to God
within the selfish desires of the old Adam (cf. \emph{Gal} 5:16;
\emph{Rom} 7:7-10). The justified also must ask God daily for
forgiveness as in the Lord\textquotesingle s Prayer (\emph{Mt}. 6:12; 1
\emph{Jn} 1:9), are ever again called to conversion and penance, and are
ever again granted forgiveness.

29.~ ~Lutherans understand this condition of the Christian as a being
"at the same time righteous and sinner." Believers are totally
righteous, in that God forgives their sins through Word and Sacrament
and grants the righteousness of Christ which they appropriate in faith.
In Christ, they are made just before God. Looking at themselves through
the law, however, they recognize that they remain also totally sinners.
Sin still lives in them (1 \emph{Jn} 1:8; \emph{Rom} 7:17,20), for they
repeatedly turn to false gods and do not love God with that undivided
love which God requires as their Creator (\emph{Deut} 6:5; \emph{Mt}
22:36-40 pr.). This contradiction to God is as such truly sin.
Nevertheless, the enslaving power of sin is broken on the basis of the
merit of Christ. It no longer is a sin that "rules" the Christian for it
is itself "ruled" by Christ with whom the justified are bound in faith.
In this life, then, Christians can in part lead a just life. Despite
sin, the Christian is no longer separated from God, because in the daily
return to baptism, the person who has been born anew by baptism and the
Holy Spirit has this sin forgiven. Thus this sin no longer brings
damnation and eternal
death.\phantomsection\label{r15}{}\hyperref[15]{{[}15{]}} Thus, when
Lutherans say that justified persons are also sinners and that their
opposition to God is truly sin, they do not deny that, despite this sin,
they are not separated from God and that this sin is a "ruled" sin. In
these affirmations, they are in agreement with Roman Catholics, despite
the difference in understanding sin in the justified.

30.~ ~Catholics hold that the grace of Jesus Christ imparted in baptism
takes away all that is sin "in the proper sense" and that is "worthy of
damnation" (\emph{Rom}
8:1).\phantomsection\label{r16}{}\href{http://www.vatican.va/roman_curia/pontifical_councils/chrstuni/documents/rc_pc_chrstuni_doc_31101999_cath-luth-joint-declaration_en.html\#_ftn16}{{[}1}\href{http://www.vatican.va/roman_curia/pontifical_councils/chrstuni/documents/rc_pc_chrstuni_doc_31101999_cath-luth-joint-declaration_en.html\#_ftn16}{6}\href{http://www.vatican.va/roman_curia/pontifical_councils/chrstuni/documents/rc_pc_chrstuni_doc_31101999_cath-luth-joint-declaration_en.html\#_ftn16}{{]}}
\textsuperscript{} There does, however, remain in the person an
inclination (concupiscence) which comes from sin and presses toward sin.
Since, according to Catholic conviction, human sins always involve a
personal element and since this element is lacking in this inclination,
Catholics do not see this inclination as sin in an authentic sense. They
do not thereby deny that this inclination does not correspond to
God\textquotesingle s original design for humanity and that it is
objectively in contradiction to God and remains one\textquotesingle s
enemy in lifelong struggle. Grateful for deliverance by Christ, they
underscore that this inclination in contradiction to God does not merit
the punishment of eternal
death\phantomsection\label{r17}{}\hyperref[17]{{[}17{]}} and does not
separate the justified person from God. But when individuals voluntarily
separate themselves from God, it is not enough to return to observing
the commandments, for they must receive pardon and peace in the
Sacrament of Reconciliation through the word of forgiveness imparted to
them in virtue of God\textquotesingle s reconciling work in Christ.
{[}See Sources for section 4.4{]}.

~

\paragraph{\texorpdfstring{\emph{4.5 Law and
Gospel}}{4.5 Law and Gospel}}\label{law-and-gospel}

31.~ ~We confess together that persons are justified by faith in the
gospel "apart from works prescribed by the law" (\emph{Rom} 3:28).
Christ has fulfilled the law and by his death and resurrection has
overcome it as a way to salvation. We also confess that
God\textquotesingle s commandments retain their validity for the
justified and that Christ has by his teaching and example expressed
God\textquotesingle s will which is a standard for the conduct of the
justified also.

32.~ ~Lutherans state that the distinction and right ordering of law and
gospel is essential for the understanding of justification. In its
theological use, the law is demand and accusation. Throughout their
lives, all persons, Christians also, in that they are sinners, stand
under this accusation which uncovers their sin so that, in faith in the
gospel, they will turn unreservedly to the mercy of God in Christ, which
alone justifies them.

33.~ ~Because the law as a way to salvation has been fulfilled and
overcome through the gospel, Catholics can say that Christ is not a
lawgiver in the manner of Moses. When Catholics emphasize that the
righteous are bound to observe God\textquotesingle s commandments, they
do not thereby deny that through Jesus Christ God has mercifully
promised to his children the grace of eternal
life.\phantomsection\label{r18}{}\hyperref[18]{{[}18{]}} {[}See Sources
for section 4.5{]}.

~

\paragraph{\texorpdfstring{\emph{4.6 Assurance of
Salvation}}{4.6 Assurance of Salvation}}\label{assurance-of-salvation}

34.~ ~We confess together that the faithful can rely on the mercy and
promises of God. In spite of their own weakness and the manifold threats
to their faith, on the strength of Christ\textquotesingle s death and
resurrection they can build on the effective promise of
God\textquotesingle s grace in Word and Sacrament and so be sure of this
grace.

35.~ ~This was emphasized in a particular way by the Reformers: in the
midst of temptation, believers should not look to themselves but look
solely to Christ and trust only him. In trust in God\textquotesingle s
promise they are assured of their salvation, but are never secure
looking at themselves.

36.~ ~Catholics can share the concern of the Reformers to ground faith
in the objective reality of Christ\textquotesingle s promise, to look
away from one\textquotesingle s own experience, and to trust in
Christ\textquotesingle s forgiving word alone (cf. \emph{Mt} 16:19;
18:18). With the Second Vatican Council, Catholics state: to have faith
is to entrust oneself totally to
God,\phantomsection\label{r19}{}\hyperref[19]{{[}19{]}} who liberates us
from the darkness of sin and death and awakens us to eternal
life.\phantomsection\label{r20}{}\hyperref[20]{{[}20{]}} In this sense,
one cannot believe in God and at the same time consider the divine
promise untrustworthy. No one may doubt God\textquotesingle s mercy and
Christ\textquotesingle s merit. Every person, however, may be concerned
about his salvation when he looks upon his own weaknesses and
shortcomings. Recognizing his own failures, however, the believer may
yet be certain that God intends his salvation. {[}See Sources for
section 4.6{]}.

~

\paragraph{\texorpdfstring{\emph{4.7 The Good Works of the
Justified}}{4.7 The Good Works of the Justified}}\label{the-good-works-of-the-justified}

37.~ ~We confess together that good works - a Christian life lived in
faith, hope and love - follow justification and are its fruits. When the
justified live in Christ and act in the grace they receive, they bring
forth, in biblical terms, good fruit. Since Christians struggle against
sin their entire lives, this consequence of justification is also for
them an obligation they must fulfill. Thus both Jesus and the apostolic
Scriptures admonish Christians to bring forth the works of love.

38.~ ~According to Catholic understanding, good works, made possible by
grace and the working of the Holy Spirit, contribute to growth in grace,
so that the righteousness that comes from God is preserved and communion
with Christ is deepened. When Catholics affirm the "meritorious"
character of good works, they wish to say that, according to the
biblical witness, a reward in heaven is promised to these works. Their
intention is to emphasize the responsibility of persons for their
actions, not to contest the character of those works as gifts, or far
less to deny that justification always remains the unmerited gift of
grace.

39.~ ~The concept of a preservation of grace and a growth in grace and
faith is also held by Lutherans. They do emphasize that righteousness as
acceptance by God and sharing in the righteousness of Christ is always
complete. At the same time, they state that there can be growth in its
effects in Christian living. When they view the good works of Christians
as the fruits and signs of justification and not as
one\textquotesingle s own "merits", they nevertheless also understand
eternal life in accord with the New Testament as unmerited "reward" in
the sense of the fulfillment of God\textquotesingle s promise to the
believer. {[}See Sources for section 4.7{]}.

~

\subsubsection{5. The Significance and Scope of the Consensus
Reached}\label{the-significance-and-scope-of-the-consensus-reached}

40.~ ~The understanding of the doctrine of justification set forth in
this Declaration shows that a consensus in basic truths of the doctrine
of justification exists between Lutherans and Catholics. In light of
this consensus the remaining differences of language, theological
elaboration, and emphasis in the understanding of justification
described in paras. 18 to 39 are acceptable. Therefore the Lutheran and
the Catholic explications of justification are in their difference open
to one another and do not destroy the consensus regarding the basic
truths.

41.~ ~Thus the doctrinal condemnations of the 16th century, in so far as
they relate to the doctrine of justification, appear in a new light: The
teaching of the Lutheran churches presented in this Declaration does not
fall under the condemnations from the Council of Trent. The
condemnations in the Lutheran Confessions do not apply to the teaching
of the Roman Catholic Church presented in this Declaration.

42.~ ~Nothing is thereby taken away from the seriousness of the
condemnations related to the doctrine of justification. Some were not
simply pointless. They remain for us "salutary warnings" to which we
must attend in our teaching and
practice.\phantomsection\label{r21}{}\hyperref[21]{{[}21{]}}

43.~ ~Our consensus in basic truths of the doctrine of justification
must come to influence the life and teachings of our churches. Here it
must prove itself. In this respect, there are still questions of varying
importance which need further clarification. These include, among other
topics, the relationship between the Word of God and church doctrine, as
well as ecclesiology, ecclesial authority, church unity, ministry, the
sacraments, and the relation between justification and social ethics. We
are convinced that the consensus we have reached offers a solid basis
for this clarification. The Lutheran churches and the Roman Catholic
Church will continue to strive together to deepen this common
understanding of justification and to make it bear fruit in the life and
teaching of the churches.

44.~ ~We give thanks to the Lord for this decisive step forward on the
way to overcoming the division of the church. We ask the Holy Spirit to
lead us further toward that visible unity which is
Christ\textquotesingle s will.

~

\section{APPENDIX}\label{appendix}

\subsubsection{\texorpdfstring{\emph{Resources for the Joint Declaration
on the Doctrine of
Justification}}{Resources for the Joint Declaration on the Doctrine of Justification}}\label{resources-for-the-joint-declaration-on-the-doctrine-of-justification}

In parts 3 and 4 of the "Joint Declaration" formulations from different
Lutheran-Catholic dialogues\\
are referred to. They are the following documents:

"All Under One Christ," Statement on the Augsburg Confession by the
Roman Catholic/Lutheran Joint Commission, 1980, in: Growth in Agreement,
edited by Harding Meyer and Lukas Vischer, New York/Ramsey, Geneva,
1984, 241-247.

Denzinger-Schönmetzer, Enchiridion symbolorum ...32nd to 36th edition
(hereafter: DS).

Denzinger-Hünermann, Enchiridion symbolorum ...since the 37th edition
(hereafter: DH).

Evaluation of the Pontifical Council for Promoting Christian Unity of
the Study Lehrverurteilungen - kirchentrennend?, Vatican, 1992,
unpublished document (hereafter: PCPCU).

Justification by Faith, Lutherans and Catholics in Dialogue VII,
Minneapolis, 1985 (hereafter: USA).

Position Paper of the Joint Committee of the United Evangelical Lutheran
Church of Germany and the LWF German National Committee regarding the
document "The Condemnations of the Reformation Era.Do They Still
Divide?" in: Lehrverurteilungen im Gespräch, Göttingen, 1993 (hereafter:
VELKD).

The Condemnations of the Reformation Era. Do they Still Divide? Edited
by Karl Lehmann and Wolfhart Pannenberg, Minneapolis, 1990 (hereafter:
LV:E)

\textbf{For 3:}The Common Understanding of Justification (paras 17 and
18) (LV:E 68f; VELKD 95)

- "... a faith centered and forensically conceived picture of
justification is of major importance for Paul and, in a sense, for the
Bible as a whole, although it is by no means the only biblical or
Pauline way of representing God\textquotesingle s saving work" (USA, no.
146).

- "Catholics as well as Lutherans can acknowledge the need to test the
practices, structures, and theologies of the church by the extent to
which they help or hinder \textquotesingle the proclamation of
God\textquotesingle s free and merciful promises in Christ Jesus which
can be rightly received only through faith\textquotesingle{} (para. 28)"
(USA, no. 153).

Regarding the "fundamental affirmation" (USA, no. 157; cf. 4) it is
said:

- "This affirmation, like the Reformation doctrine of justification by
faith alone, serves as a criterion for judging all church practices,
structures, and traditions precisely because its counterpart is
\textquotesingle Christ alone\textquotesingle{} (solus Christus). He
alone is to be ultimately trusted as the one mediator through whom God
in the Holy Spirit pours out his saving gifts. All of us in this
dialogue affirm that all Christian teachings, practices, and offices
should so function as to foster \textquotesingle the obedience of
faith\textquotesingle{} (\emph{Rom}. 1:5) in God\textquotesingle s
saving action in Christ Jesus alone through the Holy Spirit, for the
salvation of the faithful and the praise and honor of the heavenly
Father" (USA, no. 160).

- "For that reason, the doctrine of justification - and, above all, its
biblical foundation - will always retain a special function in the
church. That function is continually to remind Christians that we
sinners live solely from the forgiving love of God, which we merely
allow to be bestowed on us, but which we in no way - in however modified
a form - \textquotesingle earn\textquotesingle{} or are able to tie down
to any preconditions or postconditions. The doctrine of justification
therefore becomes the touchstone for testing at all times whether a
particular interpretation of our relationship to God can claim the name
of \textquotesingle Christian.\textquotesingle{} At the same time, it
becomes the touchstone for the church, for testing at all times whether
its proclamation and its praxis correspond to what has been given to it
by its Lord" (LV:E 69).

- "An agreement on the fact that the doctrine of justification is
significant not only as one doctrinal component within the whole of our
church\textquotesingle s teaching, but also as the touchstone for
testing the whole doctrine and practice of our churches, is - from a
Lutheran point of view - fundamental progress in the ecumenical dialogue
between our churches. It cannot be welcomed enough" (VELKD 95, 20-26;
cf. 157).

- "For Lutherans and Catholics, the doctrine of justification has a
different status in the hierarchy of truth; but both sides agree that
the doctrine of justification has its specific function in the fact that
it is \textquotesingle the touchstone for testing at all times whether a
particular interpretation of our relationship to God can claim the name
of "Christian". At the same time it becomes the touchstone for the
church, for testing at all times whether its proclamation and its praxis
correspond to what has been given to it by its Lord\textquotesingle{}
(LV:E 69). The criteriological significance of the doctrine of
justification for sacramentology, ecclesiology and ethical teachings
still deserves to be studied further" (PCPCU 96).

\textbf{For 4.1:}Human Powerlessness and Sin in Relation to
Justification (paras 19-21) (LV:E 42ff; 46; VELKD 77-81; 83f)

- "Those in whom sin reigns can do nothing to merit justification, which
is the free gift of God\textquotesingle s grace. Even the beginnings of
justification, for example, repentance, prayer for grace, and desire for
forgiveness, must be God\textquotesingle s work in us" (USA, no. 156.3).

- "\emph{Both} are concerned to make it clear that ... human beings
cannot ... cast a sideways glance at their own endeavors ... But a
response is not a \textquotesingle work.\textquotesingle{} The response
of faith is itself brought about through the uncoercible word of promise
which comes to human beings from outside themselves. There can be
\textquotesingle{}\emph{co}operation\textquotesingle{} only in the sense
that in faith the heart is involved, when the Word touches it and
creates faith" (LV:E 46f).

- "Where, however, Lutheran teaching construes the relation of God to
his human creatures in justification with such emphasis on the divine
\textquotesingle monergism\textquotesingle{} or the sole efficacy of
Christ in such a way, that the person\textquotesingle s willing
acceptance of God\textquotesingle s grace - which is itself a gift of
God - has no essential role in justification, then the Tridentine canons
4, 5, 6 and 9 still constitute a notable doctrinal difference on
justification" (PCPCU 22).

-"The strict emphasis on the passivity of human beings concerning their
justification never meant, on the Lutheran side, to contest the full
personal participation in believing; rather it meant to exclude any
cooperation in the event of justification itself. Justification is the
work of Christ alone, the work of grace alone" (VELKD 84,3-8).

\textbf{For 4.2:}Justification as Forgiveness of Sins and Making
Righteous (paras. 22-24) (USA, nos. 98-101; LV:E 47ff; VELKD 84ff; cf.
also the quotations for 4.3)

- "By justification we are both declared and made righteous.
Justification, therefore, is not a legal fiction. God, in justifying,
effects what he promises; he forgives sin and makes us truly righteous"
(USA, no. 156,5).

- "Protestant theology does not overlook what Catholic doctrine
stresses: the creative and renewing character of God\textquotesingle s
love; nor does it maintain ..God\textquotesingle s impotence toward a
sin which is \textquotesingle merely\textquotesingle{} forgiven in
justification but which is not truly abolished in its power to divide
the sinner from God" (LV:E 49).

- "The Lutheran doctrine has never understood the
\textquotesingle crediting of Christ\textquotesingle s
justification\textquotesingle{} as without effect on the life of the
faithful, because Christ\textquotesingle s word achieves what it
promises. Accordingly the Lutheran doctrine understands grace as
God\textquotesingle s favor, but nevertheless as effective power
..\textquotesingle for where there is forgiveness of sins, there is also
life and salvation\textquotesingle" (VELKD 86,15-23).

- "Catholic doctrine does not overlook what Protestant theology
stresses: the personal character of grace, and its link with the Word;
nor does it maintain ..grace as an objective
\textquotesingle possession\textquotesingle{} (even if a conferred
possession) on the part of the human being - something over which he can
dispose" (LV:E 49).

\textbf{For 4.3:}Justification by Faith and through Grace (paras.25-27)
(USA, nos. 105ff; LV:E 49-53; VELKD 87-90)

- "If we translate from one language to another, then Protestant talk
about justification through faith corresponds to Catholic talk about
justification through grace; and on the other hand, Protestant doctrine
understands substantially under the one word
\textquotesingle faith\textquotesingle{} what Catholic doctrine
(following 1 \emph{Cor}. 13:13) sums up in the triad of
\textquotesingle faith, hope, and love\textquotesingle" (LV:E 52).

- "We emphasize that faith in the sense of the first commandment always
means love to God and hope in him and is expressed in the love to the
neighbour" (VELKD 89,8-11).

- "Catholics ..teach as do Lutherans, that nothing prior to the free
gift of faith merits justification and that all of God\textquotesingle s
saving gifts come through Christ alone" (USA, no. 105).

- "The Reformers ..understood faith as the forgiveness and fellowship
with Christ effected by the word of promise itself .. This is the ground
for the new being, through which the flesh is dead to sin and the new
man or woman in Christ has life (sola fide per Christum). But even if
this faith necessarily makes the human being new, the Christian builds
his confidence, not on his own new life, but solely on
God\textquotesingle s gracious promise. Acceptance in Christ is
sufficient, if \textquotesingle faith\textquotesingle{} is understood as
\textquotesingle trust in the promise\textquotesingle{} (fides
promissionis)" (LV:E 50).

- Cf. The Council of Trent, Session 6, Chap. 7: "Consequently, in the
process of justification, together with the forgiveness of sins a person
receives, through Jesus Christ into whom he is grafted, all these
infused at the same time: faith, hope and charity" (DH 1530).

- "According to Protestant interpretation, the faith that clings
unconditionally to God\textquotesingle s promise in Word and Sacrament
is sufficient for righteousness before God, so that the renewal of the
human being, without which there can be no faith, does not in itself
make any contribution to justification" (LV:E 52).

- "As Lutherans we maintain the distinction between justification and
sanctification, of faith and works, which however implies no separation"
(VELKD 89,6-8).

- "Catholic doctrine knows itself to be at one with the Protestant
concern in emphasizing that the renewal of the human being does not
\textquotesingle contribute\textquotesingle{} to justification, and is
certainly not a contribution to which he could make any appeal before
God. Nevertheless it feels compelled to stress the renewal of the human
being through justifying grace, for the sake of acknowledging
God\textquotesingle s newly creating power; although this renewal in
faith, hope, and love is certainly nothing but a response to
God\textquotesingle s unfathomable grace" (LV:E 52f).

- "Insofar as the Catholic doctrine stresses that grace is personal and
linked with the Word, that renewal ..is certainly nothing but a response
effected by God\textquotesingle s word itself, and that the renewal of
the human being does not contribute to justification, and is certainly
not a contribution to which a person could make any appeal before God,
our objection ..no longer applies" (VELKD 89,12-21).

\textbf{For 4.4:}The Justified as Sinner (paras.28-30) (USA, nos. 102ff;
LV:E 44ff; VELKD 81ff)

- "For however just and holy, they fall from time to time into the sins
that are those of daily existence.

What is more, the Spirit\textquotesingle s action does not exempt
believers from the lifelong struggle against sinful tendencies.
Concupiscence and other effects of original and personal sin, according
to Catholic doctrine, remain in the justified, who therefore must pray
daily to God for forgiveness" (USA, no. 102).

- "The doctrines laid down at Trent and by the Reformers are at one in
maintaining that original sin, and also the concupiscence that remains,
are in contradiction to God ..object of the lifelong struggle against
sin ..{[}A{]}fter baptism, concupiscence in the person justified no
longer cuts that person off from God; in Tridentine language, it is
\textquotesingle no longer sin in the real sense\textquotesingle; in
Lutheran phraseology, it is peccatum regnatum,
\textquotesingle controlled sin\textquotesingle" (LV:E 46).

- "The question is how to speak of sin with regard to the justified
without limiting the reality of salvation. While Lutherans express this
tension with the term \textquotesingle controlled sin\textquotesingle{}
(peccatum regnatum) which expresses the teaching of the Christian as
\textquotesingle being justified and sinner at the same
time\textquotesingle{} (simul iustus et peccator), Roman Catholics think
the reality of salvation can only be maintained by denying the sinful
character of concupiscence. With regard to this question a considerable
rapprochement is reached if LV:E calls the concupiscence that remains in
the justified a \textquotesingle contradiction to God\textquotesingle{}
and thus qualifies it as sin" (VELKD 82,29-39).

\textbf{For 4.5:}Law and Gospel (paras. 31-33)

- According to Pauline teaching this topic concerns the Jewish law as
means of salvation. This law was fulfilled and overcome in Christ. This
statement and the consequences from it have to be understood on this
basis.

- With reference to Canons 19f of the Council of Trent, the VELKD
(89,28-36) says as follows:

"The ten commandments of course apply to Christians as stated in many
places of the confessions.. If Canon 20 stresses that a person ..is
bound to keep the commandments of God, this canon does not strike to us;
if however Canon 20 affirms that faith has salvific power only on
condition of keeping the commandments this applies to us. Concerning the
reference of the Canon regarding the commandments of the church, there
is no difference between us if these commandments are only expressions
of the commandments of God; otherwise it would apply to us."

- The last paragraph is related factually to 4.3, but emphasizes the
\textquotesingle convicting function\textquotesingle{} of the law which
is important to Lutheran thinking.

\textbf{For 4.6:}Assurance of Salvation (paras.34-36) (LV:E 53-56; VELKD
90ff)

- "The question is: How can, and how may, human beings live before God
in spite of their weakness, and with that weakness?" (LV:E 53).

- "The foundation and the point of departure {[}of the Reformers is{]}
..the reliability and sufficiency of God\textquotesingle s promise, and
the power of Christ\textquotesingle s death and resurrection; human
weakness, and the threat to faith and salvation which that involves"
(LV:E 56).

- The Council of Trent also emphasizes that "it is necessary to believe
that sins are not forgiven, nor have they ever been forgiven, save
freely by the divine mercy on account of Christ;" and that we must not
doubt "the mercy of God, the merit of Christ and the power and efficacy
of the sacraments; so it is possible for anyone, while he regards
himself and his own weakness and lack of dispositions, to be anxious and
fearful about his own state of grace" (Council of Trent, Session 6,
chapter 9, DH 1534).

- "Luther and his followers go a step farther. They urge that the
uncertainty should not merely be endured. We should avert our eyes from
it and take seriously, practically, and personally the objective
efficacy of the absolution pronounced in the sacrament of penance, which
comes \textquotesingle from outside.\textquotesingle{} ..Since Jesus
said, \textquotesingle Whatever you loose on earth shall be loosed in
heaven\textquotesingle{} (\emph{Matt}. 16:19), the believer ..would
declare Christ to be a liar ..if he did not rely with a rock-like
assurance on the forgiveness of God uttered in the absolution ..This
reliance can itself be subjectively uncertain - that the assurance of
forgiveness is not a security of forgiveness (securitas); but this must
not be turned into yet another problem, so to speak: the believer should
turn his eyes away from it, and should look only to
Christ\textquotesingle s word of forgiveness" (LV:E 53f).

- "Today Catholics can appreciate the Reformer\textquotesingle s efforts
to ground faith in the objective reality of Christ\textquotesingle s
promise, \textquotesingle whatsoever you loose on earth
....\textquotesingle{} and to focus believers on the specific word of
absolution from sins. ..Luther\textquotesingle s original concern to
teach people to look away from their experience, and to rely on Christ
alone and his word of forgiveness {[}is not to be condemned{]}" (PCPCU
24).

- A mutual condemnation regarding the understanding of the assurance of
salvation "can even less provide grounds for mutual objection today -
particularly if we start from the foundation of a biblically renewed
concept of faith. For a person can certainly lose or renounce faith, and
self-commitment to God and his word of promise. But if he believes in
this sense, he \emph{cannot at the same time} believe that God is
unreliable in his word of promise. In this sense it is true today also
that - in Luther\textquotesingle s words - faith \emph{is} the assurance
of salvation" (LV:E 56).

- With reference to the concept of faith of Vatican II, see Dogmatic
Constitution on Divine Revelation, no. 5: "\textquotesingle The
obedience of faith\textquotesingle{} ..must be given to God who reveals,
an obedience by which man entrusts his whole self freely to God,
offering \textquotesingle the full submission of intellect and will to
God who reveals,\textquotesingle{} and freely assenting to the truth
revealed by Him."

- "The Lutheran distinction between the certitude (certitudo) of faith
which looks alone to Christ and earthly security (securitas), which is
based on the human being, has not been dealt with clearly enough in the
LV. The question whether a Christian ``has believed fully and
completely'' (LV:E 53) does not arise for the Lutheran understanding,
since faith never reflects on itself, but depends completely on God,
whose grace is bestowed through word and sacrament, thus from outside
(extra nos)" (VELKD 92,2-9).

\textbf{For 4.7:}The Good Works of the Justified (paras.37-39) (LV:E
66ff, VELKD 90ff)

- "But the Council excludes the possibility of earning \emph{grace} -
that is, justification - (can. 2; DS 1552) and bases the earning or
merit of \emph{eternal life} on the gift of grace itself, through
membership in Christ (can. 32: DS 1582). Good works are
\textquotesingle merits\textquotesingle{} as a \emph{gift}. Although the
Reformers attack \textquotesingle Godless trust\textquotesingle{} in
one\textquotesingle s own works, the Council explicitly excludes any
notion of a claim or any false security (cap. 16: DS 1548f). It is
evident ..that the Council wishes to establish a link with Augustine,
who introduced the concept of merit, in order to express the
responsibility of human beings, in spite of the
\textquotesingle bestowed\textquotesingle{} character of good works"
(LV:E 66).

- If we understand the language of "cause" in Canon 24 in more personal
terms, as it is done in chapter 16 of the Decree on Justification, where
the idea of communion with Christ is foundational, then we can describe
the Catholic doctrine on merit as it is done in the first sentence of
the second paragraph of 4.7: growth in grace, perseverance in
righteousness received from God and a deeper communion with Christ.

- "Many antitheses could be overcome if the misleading word
\textquotesingle merit\textquotesingle{} were simply to be viewed and
thought about in connection with the true sense of the biblical term
\textquotesingle wage\textquotesingle{} or reward" (LV:E 67).

- "The Lutheran confessions stress that the justified person is
responsible not to lose the grace received but to live in it ..Thus the
confessions can speak of a preservation of grace and a growth in it. If
righteousness in Canon 24 is understood in the sense that it affects
human beings, then it does not strike to us. But if
\textquotesingle righteousness\textquotesingle{} in Canon 24 refers to
the Christian\textquotesingle s acceptance by God, it strikes to us; for
this righteousness is always perfect; compared with it the works of
Christians are only \textquotesingle fruits\textquotesingle{} and
\textquotesingle signs\textquotesingle" (VELKD 94,2-14).

- "Concerning Canon 26, we refer to the Apology where eternal life is
described as reward: \textquotesingle..We grant that eternal life is a
reward because it is something that is owed - not because of our merits
but because of the promise\textquotesingle" (VELKD 94,20-24).

~

~

~

\hyperref[r1]{}\phantomsection\label{1}{}{[}\hyperref[r1]{1}{]} The
Smalcald Articles, II,1; Book of Concord, 292.

\hyperref[r2]{}\phantomsection\label{2}{}{[}\hyperref[r2]{2}{]} "Rector
et judex super omnia genera doctrinarum" Weimar Edition of
Luther\textquotesingle s Works (WA), 39,I,205.

\hyperref[r3]{}\phantomsection\label{3}{}{[}\hyperref[r3]{3}{]} It
should be noted that some Lutheran churches include only the Augsburg
Confession and Luther\textquotesingle s Small Catechism among their
binding confessions. These texts contain no condemnations about
justification in relation to the Roman Catholic Church.

\hyperref[r4]{}\phantomsection\label{4}{}{[}\hyperref[r4]{4}{]} Report
of the Joint Lutheran-Roman Catholic Study Commission, published in
Growth in Agreement (New York; Geneva, 1984), pp. 168-189.

\hyperref[r5]{}\phantomsection\label{5}{}{[}\hyperref[r5]{5}{]}
Published by the Lutheran World Federation (Geneva, 1994).

\hyperref[r6]{}\phantomsection\label{6}{}{[}\hyperref[r6]{6}{]} Lutheran
and Catholics in Dialogue VII (Minneapolis, 1985).

\hyperref[r7]{}\phantomsection\label{7}{}{[}\hyperref[r7]{7}{]}
Minneapolis, 1990.

\hyperref[r8]{}\phantomsection\label{8}{}{[}\hyperref[r8]{8}{]}
"Gemeinsame Stellungnahme der Arnoldshainer Konferenz, der Vereinigten
Kirche und des Deutschen Nationalkomitees des Lutherischen Weltbundes
zum Dokument \textquotesingle Lehrverurteilungen -
kirchentrennend?\textquotesingle," Ökumenische Rundschau 44 (1995):
99-102; See also the position papers which underlie this resolution, in
Lehrverurteilungen im Gespräch, Die ersten offiziellen Stellungnahmen
aus den evangelischen Kirchen in Deutschland (Göttingen: Vandenhoeck \&
Ruprecht, 1993).

\hyperref[r9]{}\phantomsection\label{9}{}{[}\hyperref[r9]{9}{]} The word
"church" is used in this Declaration to reflect the self-understandings
of the participating churches, without intending to resolve all the
ecclesiological issues related to this term.\href{1}{}

\hyperref[r10]{}\phantomsection\label{10}{}{[}\hyperref[r10]{10}{]} Cf.
"Malta Report," paras. 26-30; Justification by Faith, paras. 122-147. At
the request of the US dialogue on justification, the non-Pauline New
Testament texts were addressed in Righteousness in the New Testament, by
John Reumann, with responses by Joseph A. Fitzmyer and Jerome D. Quinn
(Philadelphia; New York:1982), pp. 124-180. The results of this study
were summarized in the dialogue report Justification by Faith in paras.
139-142.

\hyperref[r11]{}\phantomsection\label{11}{}{[}\hyperref[r11]{11}{]} "All
Under One Christ," para. 14, in Growth in Agreement, 241-247.

\hyperref[r12]{}\phantomsection\label{12}{}{[}\hyperref[r12]{12}{]} Cf.
WA 8:106; American Edition 32:227.

\hyperref[r13]{}\phantomsection\label{13}{}{[}\hyperref[r13]{13}{]} Cf.
DS 1528

\hyperref[r14]{}\phantomsection\label{14}{}{[}\hyperref[r14]{14}{]} Cf.
DS 1530.

\hyperref[r15]{}\phantomsection\label{15}{}{[}\hyperref[r15]{15}{]} Cf.
Apology II:38-45; Book of Concord, 105f.

\hyperref[r16]{}\phantomsection\label{16}{}{[}\hyperref[r16]{16}{]} Cf.
DS 1515.

\hyperref[r17]{}\phantomsection\label{17}{}{[}\hyperref[r17]{17}{]} Cf.
DS 1515.

\hyperref[r18]{}\phantomsection\label{18}{}{[}\hyperref[r18]{18}{]} Cf.
DS 1545.

\hyperref[r19]{}\phantomsection\label{19}{}{[}\hyperref[r19]{19}{]} Cf.
DV 5.

\href{http://www.vatican.va/roman_curia/pontifical_councils/chrstuni/documents/rc_pc_chrstuni_doc_31101999_cath-luth-joint-declaration_en.html\#_ftnref20}{}\phantomsection\label{20}{}{[}\hyperref[r20]{20}{]}
Cf. DV 5.

\phantomsection\label{21}{}{[}\hyperref[r21]{21}{]} Condemnations of the
Reformation Era, 27.

~

\textsuperscript{{[}\emph{Information Service~}98 (1998/III) 81-89{]}}

~
